\scanpage{67}
% archive foliation: 69
et considérons ``l'espace de
Teichmüller à bord'' correspondant
\deleted{$\unsure{T_{g\nu}}$} $TB_{g\nu}$, qui est \deleted{isom.}\inserted{homéo.} à $\mathbb R^{6g-6+3\nu}$
ou mieux, à $\mathbb C^{3g-3+\nu}\times\mathbb R^\nu$. On a une
application
\diagram{\begin{tikzcd}[ampersand replacement=\&]
\underset{\simeq\mathbb C^{3g-3+\nu}\times\mathbb R^\nu}{TB_{g,\nu}} \arrow[r] \&
\underset{\mathbb C^{3g-3+\nu}\deleted{\times\mathbb R^\nu}}{T_{g\nu}}
\end{tikzcd}}
de bouchage des trous, je présume
qu'elle est fibrante \unsure{de faç.},
plus précise, que \deleted{\unsure{dans}\ \unsure{la}}\inserted{c'est la projection}
canonique. \deleted{de} En effet, cela doit
provenir du fait que pour une
sphère à bord \inserted{\unsure{déformante|déformable}} de type $(0,3)$, \deleted{dont}
quand le modèle combinatoire
(\unsure{triangle conducteur}) est fixé, et
qu'on varie seulement \deleted{les} un, deux
ou trois des paramètres ``longueurs''
correspondant aux trois trous, la
sphère ``à bord et ponctuée'' (mono, bi ou
triponctuée suivant les cas) \deleted{\unsure{et le}}\inserted{\unsure{associée|acquise}}
